\section{The auxiliary-prime torsion sum for a nonstandard lattice}

% BEGIN A76V1
\paragraph{A76V1. The primary kernel and its fixed subgroup.}
Suppose $\ell\mid N$, $\ell\nmid M$. Keep an original summand
$X_q$ of (ICL.13). Work over the chart of R76V1, allowing
further finite extensions with the same normalization of $\nu$.
Write $S_\ell=(\mathcal O_E\otimes\mathbb Z_\ell)^n$,
$\Lambda_\ell=\Lambda\otimes\mathbb Z_\ell$ and $r=D_w.p+D_w.q$.
Choose a rational unitary matrix $\gamma$ with
$\gamma\Lambda_\ell=S_\ell$ and $\gamma$ integral with integral
inverse over 2. Here $\ell$ is odd. At split primes the choice
of one local basis determines its dual on the paired prime. At
nonsplit primes the host is unramified over its maximal real
subfield; the norm-one orthogonal basis of an integral unimodular
lattice is obtained by splitting off unit vectors as in S74B4.
At 2 the standard lattice already agrees with $\Lambda$ since
$N$ is odd. Weak approximation gives the rational element:
the conditions on the local matrices are open, and on the
unitary Cayley chart approximation is approximation of
skew-adjoint matrices over the real subfield. The chart is
dense at each of these places. At $\ell$, $NS_\ell\subseteq
\Lambda_\ell$ gives $N\gamma S_\ell\subseteq S_\ell$. Choose
$k>0$ prime to $2\ell$ clearing the remaining denominators of
$N\gamma$, and put $L=kN$, $U=L\gamma$. Then $U$ is an
integral matrix with $U^*U=UU^*=L^2 I_n$.

Let $A$ denote the generic abelian variety at the framed test.
Use its geometric $\ell$-power torsion to describe the primary
subgroup $W$ of the isotropic tuples in R76V1. The Tate lattice
$T_\ell A$ is free of rank $r$ over
$\mathcal O_E\otimes\mathbb Z_\ell$, by the polarized generic
marking used in the count T73C1. Equivalently choose a basis
of that lattice at each prime over $\ell$. In each of its $r$
copies the primary points lie in
$(E\otimes\mathbb Q_\ell)^n/S_\ell$. The matrix tests defining
$\mathrm{InZh}$ say that $N$ times a representative pairs
integrally with every row of the ordinary dual of $\Lambda_\ell$.
This is precisely membership in $\Lambda_\ell$ by biduality.
Indeed those integral-matrix tests localize: the rows in the
global ordinary dual have denominators dividing $N$ because
$NS\subseteq\Lambda$, and they span the local dual. Therefore
\[
 W=(N^{-1}\Lambda_\ell/S_\ell)^r=\ker(U)_\ell,\qquad
 W_0=W\cap A[N]^n
      =\big(N^{-1}(\Lambda_\ell\cap S_\ell)/S_\ell\big)^r .
\]
Multiplication by the unit $k$ preserves the localized lattice in
this equality. In particular $|W_0|$ is exactly the $\ell$-part
of $\mathrm{mIdx}=[\Lambda\cap S:NS]^r$.

Fix also the prime-to-$\ell$ component $z'$ of the enumerated
tuple. The other translate $t_q$ has order prime to $\ell$
since it is killed by a power of the half ideal, hence by a
power of $M$. Start with the $[N]$-pullback of the frame section,
apply these two prime-to-$\ell$ translations and pull the columns
back by $[k]$. Their tensor product on $A^n$ is denoted by $F$.
Write $F_a$ for its translate by the product of the symmetric
lifts at the coordinates of $a\in W$, at tensor power $L^2$.
Multiplication by $k$ permutes $W$. Pullback along $[k]$
intertwines these lifts with those at $[k]a$ at level $N^2$.
The prime-to-$\ell$ lifts commute with the primary ones because
their commutators have coprime annihilators. Thus these are,
up to this permutation, exactly the normalized $[k]$-pullbacks
of the products in the $z'$ part of $X_q$. Each has an integral
germ by reachability and R76V2.

On $W$ the tensor product lifts act as a group. Indeed the
product commutator is one by isotropy. The possible product
scalar for composing two symmetric lifts has square equal to
the product commutator and has odd order, so equals one.
For $a\in W_0$ the lift at $[k]a$ acts trivially on the original
$[N]$-pullback by its deck action, and commutes with the two
prime-to-$\ell$ translations. Consequently $F_{a+w}=F_a$
for $w\in W_0$. We can divide the sum by $|W_0|$ before taking
any jets: $F_{\rm sum}=\sum_{a\in W/W_0}F_a$ is an integral
germ and the action of each $w\in W$ on it is trivial.
% END A76V1

% BEGIN A76V2
\paragraph{A76V2. Integral descent on the formal product.}
Write $\mathcal L$ for the external product of the $n$
rigidified theta lines on the integral semiabelian model.
There is a rigidified comparison
$U^*\mathcal L\simeq\mathcal L^{L^2}$. To construct it, write
the line in each factor as the primitive diagonal pullback of
the biextension at $(x,[w_0]x)$. Expand the two occurrences
of the matrix $U$ using the two addition laws of the
biextension. The factor at $(x_j,x_m)$ is obtained by applying
$\sum_i\overline{U}_{ij}U_{im}w_0$ to the second entry, by the
integral adjoint isomorphisms. The relation $U^*U=L^2 I$
cancels the off-diagonal factors and leaves the indicated
power on each diagonal. The cubical associativity and the
rigidifications on the axes give the identity on the origin
fiber. The same construction commutes with inversion and gives
an integral unit on the formal product for any choice of its
compatible generators. In the canonical gauge, write
$Q(x)=\sum_j h(x_j,[w_0]x_j)$. The scalar for expressing the
pullback of a target germ in the source trivialization is
\[
               \kappa(x)=\exp\big(Q(Ux)-L^2 Q(x)\big).
\]
Indeed the two addition and adjoint defects in the preceding
expansion are exactly the logarithms of the integral transition
units defining this comparison. The mixed terms cancel by the
same matrix equality. The unit is normalized by $\kappa(0)=1$.
This also shows directly that no formal logarithm is evaluated
on a nonzero torsion point in constructing the comparison.

For $a\in\ker(U)_\ell$ the ordinary deck action on $U^*\mathcal L$
is the tensor product symmetric lift of A76V1. Over the generic
abelian product both are isomorphisms above translation by $a$,
so their ratio is a scalar. Both have odd order and are taken
to their inverses by the rigidified inversion isomorphism.
The ratio therefore has square one and odd order, hence is one.
Use this generic identity for the algebraic sections of A76V1.
In particular the integral germ $f=\kappa^{-1}F_{\rm sum}$
is invariant under ordinary translation by every point of
$\ker(U)$ in the formal group.

Here is the integral descent argument, including the behavior at
a semiabelian special fiber. In analytic coordinates of the
formal completion write its complete local ring as
$R=V'[[Y_1,\ldots,Y_m]]$ and the source ring as
$S'=V'[[X_1,\ldots,X_m]]$, where $m=n\dim A$.
The substitution $U^*:R\to S'$ is finite flat. In fact the
closed-fiber kernel of $U$ at the origin is zero-dimensional:
it is killed by $[L^2]$ by the adjoint identity, and on a
semiabelian variety multiplication has finite kernel, including
on the torus. The images of the $Y_i$ together with the
uniformizer thus generate a primary ideal for the maximal ideal.
Finiteness follows by lifting a finite basis modulo that ideal
and completeness. These $m+1$ parameters are a regular sequence
in the regular local source, so the finite module is free over
the regular local target. Its rank $d_0$ is the degree of the
finite flat formal kernel, obtained by setting $Y=0$. In
characteristic zero it is \'etale: in invariant coordinates the
linear map of $U$ has inverse $U^*/L^2$ over the fraction field.
After finite splitting, there are $d_0$ kernel points, all
specializing to the origin. They are genuine geometric points
of the original kernel by the formal chart (coordinates in the
maximal ideal converge over the complete trait). They have
$\ell$-power order: every integer prime to $\ell$ induces an
automorphism on the formal group. Hence all lie in $W$, also
after any additional splitting extension with renormalized $e'$.

Translation by these points gives $d_0$ distinct automorphisms
of $S'$ over $R$; distinct images of the origin distinguish
them. They extend to $\operatorname{Frac}(S')$, of degree
$d_0$ over $\operatorname{Frac}(R)$. Their fixed field is
therefore $\operatorname{Frac}(R)$. The invariant $f\in S'$
lies in this fixed field, is integral over $R$ by finiteness,
and $R$ is integrally closed. It follows that $f=U^*g$ for
some $g\in R$. We have proved
$F_{\rm sum}=\kappa\,U^*g$ with $g$ an integral joint germ.
This proof uses only the fixed-subgroup cardinal already
removed in A76V1.
% END A76V2

% BEGIN A76V3
\paragraph{A76V3. Pushforward of the polynomial and the sharp auxiliary bound.}
In the harmonic evaluation of the joint germ of A76V2 the
original connection multiplies by $\exp(L^2Q(x))$. In the
display of A76V2 this cancels the second term in the exponent
of $\kappa$, leaving $U^*(g\exp Q)$. The invariant-direction
chain rule takes each column derivative to the sum of the
target column derivatives with coefficients from that column of
the matrix $U$ under the relevant embedding. Thus the new
polynomial is exactly $\mathrm{push}_U P$. Evaluation on the
$[k]$-pullbacks in A76V1 contributes $k^d$ compared with
the original $z'$ sum, since $[k]$ multiplies each direction
by $k$ and $d=\deg P$. For $P\ne0$ the original factor
$N^{-d}$ together with this pullback factor is $L^{-d}$.
Homogeneity yields $L^{-d}\mathrm{push}_U P
=\mathrm{push}_\gamma P$. The last polynomial has integral
coefficients at the prescribed test: the generators
$\mathrm{letB}(\sigma,i,v)$ for $v\in\Lambda$ are sent to
the corresponding letters for $\gamma v$, integral at every
prime above $\ell$. Unitarity also preserves harmonicity
and row degrees.

More explicitly, the original contribution of $z'$ to $X_q$ is
$|W_0|/\mathrm{mIdx}$ times the level-one evaluation of $g$
with polynomial $\mathrm{push}_\gamma P$. The prefactor is a
unit by A76V1. P76V2 applies to any integral joint germ here:
in its labelled product rule, the derivatives hitting $g$
remain integral even when they hit different copies. For the
exponential, $Q$ is the sum of the individual $q$'s, so a
mixed derivative involving different copies vanishes. As
before, only disjoint pairs can contribute to the denominator.
This gives valuation at least $-\lfloor d/2\rfloor c_\ell$
on this coset contribution. Sum over $z'$ and use the
valuation inequality. The full $X_q$ satisfies (ICL.14)
with $\alpha=1$, including $P=0$. The descent through any
additional finite splitting preserves the normalization
$\mathrm{ord}_{V'}/(be)$.

Together with P76V3 this supplies (ICL.13)--(ICL.15) on the
original identity cover with squarefree $M$, with the
specified ramified loss and the actual finite averages.
ICL17--ICL18 now give the uniform scalar of (ICL.17) on
that cover, independent of the auxiliary level and polynomial
once the displayed row data and $\beta_H(\Lambda)$ are fixed.
% END A76V3
